\section*{Chapter \ref{ch:b1:series} --- Numerical Series}

\begin{solution}{exo:b1:series:1}
$\dfrac{1}{n(n+1)} = \dfrac1n - \dfrac{1}{n+1}$: telescoping, $S_N =
1 - \frac{1}{N+1} \to 1$. Convergent, sum $1$.

$\ln\bigl(1 - \frac{1}{n^2}\bigr) = \ln\frac{(n-1)(n+1)}{n^2} =
\ln\frac{n-1}{n} - \ln\frac{n}{n+1}$: telescoping again, $S_N =
\ln\frac12 - \ln\frac{N}{N+1} \to -\ln 2$. Convergent, sum $-\ln 2$.

$\dfrac{3^n + 4^n}{5^n} = \bigl(\frac35\bigr)^n +
\bigl(\frac45\bigr)^n$: two convergent geometric series, sum
$\frac{1}{1 - 3/5} + \frac{1}{1 - 4/5} = \frac52 + 5 = \frac{15}{2}$.
\end{solution}

\begin{solution}{exo:b1:series:2}
Ratio test throughout.

$\frac{u_{n+1}}{u_n} = \frac{(n+1)^2}{2n^2} \to \frac12 < 1$:
convergent.

$\frac{u_{n+1}}{u_n} = \frac{(n+1)!\,n^n}{n!\,(n+1)^{n+1}} =
\bigl(\frac{n}{n+1}\bigr)^n = \bigl(1 + \frac1n\bigr)^{-n} \to
\frac1\eu < 1$: convergent.

With the factor $2^n$: ratio $\to \frac2\eu < 1$: convergent.

With $3^n$: ratio $\to \frac3\eu > 1$: divergent (terms tend to
$+\infty$).
\end{solution}

\begin{solution}{exo:b1:series:3}
All nonnegative terms; use equivalents
(\cref{thm:b1:series:positive}).

$\sin\frac{1}{n^2} \sim \frac{1}{n^2}$: convergent (Riemann $\alpha
= 2$).

$1 - \cos\frac1n \sim \frac{1}{2n^2}$: convergent.

$\frac{1}{\sqrt{n(n+1)}} \sim \frac1n$: divergent.

$\frac{\ln n}{n^2} = \frac{1}{n^{3/2}}\cdot\frac{\ln n}{n^{1/2}}$
and $\frac{\ln n}{\sqrt n} \to 0$
(\cref{prop:b1:functions:powerrules}): so $\frac{\ln n}{n^2} \leq
\frac{1}{n^{3/2}}$ for large $n$: convergent.
\end{solution}

\begin{solution}{exo:b1:series:4}
For $n \geq 2$: $\frac{1}{n^2} \leq \frac{1}{n(n-1)} =
\frac{1}{n-1} - \frac1n$. Hence
\[
\sum_{n=1}^{N} \frac{1}{n^2} \leq 1 + \sum_{n=2}^{N}
\Bigl(\frac{1}{n-1} - \frac1n\Bigr) = 1 + 1 - \frac1N < 2 :
\]
partial sums increasing and bounded by $2$: convergence
(\cref{thm:b1:series:positive}), sum $\leq 2$. (The exact value
$\frac{\pi^2}{6}$ is a second-year celebration.)
\end{solution}

\begin{solution}{exo:b1:series:5}
$\sum \frac{(-1)^n}{\sqrt n}$: alternating with $\frac{1}{\sqrt n}
\downarrow 0$: convergent (\cref{thm:b1:series:alternating}); not
absolutely ($\alpha = \frac12 \leq 1$).

$\sum \frac{(-1)^n}{n + (-1)^n}$: the sequence $\frac{1}{n +
(-1)^n}$ is \emph{not} decreasing ($\frac{1}{n+1}$ then
$\frac{1}{n}$ alternate badly), so the test does not apply directly.
Expand:
\[
\frac{(-1)^n}{n + (-1)^n}
= \frac{(-1)^n}{n}\cdot\frac{1}{1 + \frac{(-1)^n}{n}}
= \frac{(-1)^n}{n} - \frac{1}{n^2} +
O\Bigl(\frac{1}{n^3}\Bigr):
\]
the first series converges (alternating), $\sum \frac{1}{n^2}$
converges, the $O(n^{-3})$ converges absolutely: the sum of three
convergent series converges.

$\sin\bigl(\pi\sqrt{n^2+1}\bigr)$: write $\sqrt{n^2 + 1} = n +
\frac{1}{2n} + \varepsilon_n$ with $\varepsilon_n = O(n^{-3})$;
then, by $\pi$-periodicity of $\sin$ up to sign,
\[
\sin\bigl(\pi\sqrt{n^2+1}\bigr)
= (-1)^n \sin\Bigl(\frac{\pi}{2n} + \pi\varepsilon_n\Bigr) .
\]
Set $\theta_n = \frac{\pi}{2n} + \pi\varepsilon_n$ and $a_n =
\sin\theta_n$. For large $n$, $\theta_n \in \intoo{0}{\frac\pi2}$
and
\[
\theta_n - \theta_{n+1} = \frac{\pi}{2n(n+1)} +
\pi(\varepsilon_n - \varepsilon_{n+1})
= \frac{\pi}{2n^2} + O\Bigl(\frac{1}{n^3}\Bigr) > 0
\]
eventually, so $(\theta_n)$ decreases to $0$; since $\sin$ is
increasing on $\intcc{0}{\frac\pi2}$, $(a_n)$ decreases to $0$ as
well. The alternating test applies: convergent --- not absolutely,
since $a_n \sim \frac{\pi}{2n}$.
\end{solution}

\begin{solution}{exo:b1:series:6}
$f(t) = \frac{1}{t(\ln t)^\alpha}$ is positive, continuous,
decreasing on $\intco{2}{+\infty}$. Substituting $u = \ln t$:
\[
\int_2^x \frac{\dd t}{t(\ln t)^\alpha}
= \int_{\ln 2}^{\ln x} \frac{\dd u}{u^\alpha},
\]
bounded as $x \to \infty$ iff $\alpha > 1$
(\cref{thm:b1:series:riemann}'s computation). By integral
comparison: convergence iff $\alpha > 1$. (These \emph{Bertrand-type}
series show how fine the boundary of convergence is: $n\ln n$
diverges, $n(\ln n)^{1.01}$ converges.)
\end{solution}

\begin{solution}{exo:b1:series:7}
By the tangent-line bounds of \cref{exo:b1:derivative:3} rewritten
via expansions: for $x = \frac1n \in \intoc{0}{1}$, Taylor--Lagrange
for $\ln(1+x)$ at order $1$ gives $\ln(1 + x) = x - \frac{x^2}{2(1 +
c)^2}$ for some $c \in \intoo{0}{x}$, so
\[
0 \leq u_n = \frac1n - \ln\Bigl(1 + \frac1n\Bigr) \leq
\frac{1}{2n^2} .
\]
Comparison with the Riemann series: $\sum u_n$ converges. Its partial
sum telescopes the logarithms:
\[
\sum_{n=1}^{N} u_n = H_N - \ln(N+1)
\]
(since $\sum_{n\leq N} \ln\frac{n+1}{n} = \ln(N+1)$). So $H_N -
\ln(N+1)$ converges; adding $\ln\frac{N+1}{N} \to 0$, the sequence
$H_N - \ln N$ converges. Its limit is $\gamma \approx 0.5772$.
\end{solution}

\begin{solution}{exo:b1:series:8}
$\dfrac{1}{n(n+2)} = \dfrac{1/2}{n} - \dfrac{1/2}{n+2}$: the partial
sum telescopes with a lag of $2$,
\[
S_N = \frac12\Bigl(1 + \frac12 - \frac{1}{N+1} -
\frac{1}{N+2}\Bigr) \longrightarrow \frac34 .
\]

$\sum \frac{n}{2^n}$: for $\abs x < 1$, differentiating the finite
geometric sum and passing to the limit (all series here converge
absolutely, ratio test): from $\sum_{n\geq0} x^n = \frac{1}{1-x}$,
one gets by direct computation with partial sums
\[
\sum_{n=1}^{N} n x^{n-1}
= \frac{1 - (N+1)x^N + N x^{N+1}}{(1 - x)^2}
\xrightarrow[N\to\infty]{} \frac{1}{(1-x)^2}
\quad (\abs x < 1),
\]
(the boundary terms $N x^N \to 0$). At $x = \frac12$:
$\sum_{n\geq1} n\bigl(\frac12\bigr)^{n-1} = 4$, so $\sum_{n\geq0}
\frac{n}{2^n} = \frac12 \times 4 = 2$.
\end{solution}

\begin{solution}{exo:b1:series:9}
Group the terms of $\sum u_n$ in packets between powers of $2$.
Upper packets: for $2^k \leq n < 2^{k+1}$ there are $2^k$ terms,
each $\leq u_{2^k}$:
\[
\sum_{n=1}^{2^{K+1}-1} u_n
= \sum_{k=0}^{K} \sum_{n=2^k}^{2^{k+1}-1} u_n
\leq \sum_{k=0}^{K} 2^k u_{2^k} .
\]
Lower packets: each term of the same packet is $\geq u_{2^{k+1}}$,
so $\sum_{n=2^k}^{2^{k+1}-1} u_n \geq 2^k u_{2^{k+1}} = \frac12
\cdot 2^{k+1} u_{2^{k+1}}$, whence
\[
\sum_{n=1}^{2^{K+1}-1} u_n \geq \frac12 \sum_{k=1}^{K+1} 2^{k}
u_{2^{k}} .
\]
Both partial-sum comparisons go both ways (nonnegative terms,
\cref{thm:b1:series:positive}): the two series have the same nature.

Riemann: $u_n = n^{-\alpha}$ gives $2^k u_{2^k} =
2^{k(1-\alpha)}$, a geometric series, convergent iff $2^{1 - \alpha}
< 1$ iff $\alpha > 1$. Bertrand
(\cref{exo:b1:series:6}): $u_n = \frac{1}{n(\ln n)^\alpha}$ gives
$2^k u_{2^k} = \frac{1}{(k\ln 2)^\alpha}$, a Riemann series in $k$:
convergent iff $\alpha > 1$.
\end{solution}

\begin{solution}{exo:b1:series:10}
Finite geometric identity with ratio $-t^2$:
\[
\frac{1}{1 + t^2} = \sum_{k=0}^{n-1} (-1)^k t^{2k}
+ \frac{(-1)^n t^{2n}}{1 + t^2} .
\]
Integrate over $\intcc{0}{1}$ (the left side integrates to $\arctan
1 = \frac\pi4$, \cref{prop:b1:functions:arcderiv}):
\[
\frac{\pi}{4} = \sum_{k=0}^{n-1} \frac{(-1)^k}{2k+1}
+ (-1)^n \int_0^1 \frac{t^{2n}}{1+t^2}\,\dd t,
\qquad
0 \leq \int_0^1 \frac{t^{2n}}{1+t^2}\,\dd t \leq \int_0^1 t^{2n}\dd
t = \frac{1}{2n+1} .
\]
Letting $n \to \infty$ proves the formula, and the displayed bound
on the integral is exactly the remainder bound: after summing up to
$N$ (i.e.\ $n = N + 1$ terms), $\abs{R_N} \leq \frac{1}{2N + 3}$.
\end{solution}

\begin{solution}{exo:b1:series:11}
$n^{1/n} = \eu^{\frac{\ln n}{n}} \to \eu^0 = 1$
(\cref{prop:b1:functions:powerrules}). Hence
\[
\frac{1}{n^{1 + 1/n}} = \frac{1}{n}\,\eu^{-\frac{\ln n}{n}}
\sim \frac{1}{n} ,
\]
and the equivalents test (\cref{thm:b1:series:positive})
compares with the divergent harmonic series: \emph{divergent},
although every exponent $1 + \frac1n$ exceeds $1$. The Riemann
criterion concerns a \emph{fixed} exponent $\alpha$; an exponent
sliding down to $1$ can lose all its margin, as here.
\end{solution}

\begin{solution}{exo:b1:series:12}
Let $\varepsilon > 0$. By the Cauchy criterion for the
convergent series (\cref{thm:b1:seq:complete} applied to the
partial sums), there is $N$ with $\sum_{k=n+1}^{2n} u_k \leq
\varepsilon$ for $n \geq N$. By monotonicity each of these $n$
terms is $\geq u_{2n}$:
\[
n\,u_{2n} \leq \sum_{k=n+1}^{2n} u_k \leq \varepsilon
\quad\Longrightarrow\quad 2n\,u_{2n} \leq 2\varepsilon ,
\]
and for odd indices $(2n+1)\,u_{2n+1} \leq (2n+1)\,u_{2n} \leq
2\bigl(2n\,u_{2n}\bigr) \leq 4\varepsilon$ for $n \geq N$: in
both parities, $n u_n \to 0$.

Converse false: $u_n = \frac{1}{n\ln n}$ has $n u_n =
\frac{1}{\ln n} \to 0$, yet the series diverges
(\cref{exo:b1:series:6}, $\alpha = 1$). Monotonicity necessary:
let $u_n = \frac1n$ when $n$ is a perfect square and $u_n =
2^{-n}$ otherwise: the series converges (the square terms sum
like $\sum \frac{1}{k^2}$, the rest geometrically), but $n u_n =
1$ along the squares.
\end{solution}

\begin{solution}{pb:b1:series:1}
\textbf{1.} $a_{n+1} - a_n = \frac{1}{n+1} - \ln\frac{n+1}{n}
\leq 0$ because $\ln(1 + \frac1n) \geq \frac{1/n}{1 + 1/n} =
\frac{1}{n+1}$; and $b_{n+1} - b_n = \frac{1}{n+1} -
\ln\frac{n+2}{n+1} \geq 0$ because $\ln(1 + \frac{1}{n+1}) \leq
\frac{1}{n+1}$. Their gap $a_n - b_n = \ln(1 + \frac1n) \to 0$:
adjacent (\cref{thm:b1:seq:adjacent}), with common limit $\lim
a_n = \gamma$ (\cref{exo:b1:series:7}), and $b_n \leq \gamma
\leq a_n$.

\textbf{2.} $b_{10} = 2.928968 - \ln 11 = 0.5311$ and $a_{10} =
2.928968 - \ln 10 = 0.6264$: so $\gamma \in
\intcc{0.5311}{0.6264}$. The gap is $\ln 1.1 \approx 0.095$ and
shrinks like $\frac1n$: four decimals ($\text{gap} \leq
10^{-4}$) would need $n \approx 10^4$ --- the brackets are
correct but slow.

\textbf{3.} Telescoping $a_n - a_{m+1} = \sum_{k=n}^{m} w_k$ and
letting $m \to \infty$: $a_n - \gamma = \sum_{k\geq n} w_k$
(limit of partial sums). Moreover
\[
w_k = \int_k^{k+1} \frac{\dd t}{t} - \frac{1}{k+1}
= \int_k^{k+1} \Bigl(\frac1t - \frac{1}{k+1}\Bigr)\dd t
= \int_k^{k+1} \frac{k + 1 - t}{t\,(k+1)}\,\dd t .
\]
On $\intcc{k}{k+1}$: $\frac{1}{(k+1)^2} \leq \frac{1}{t(k+1)}
\leq \frac{1}{k(k+1)}$, and $\int_k^{k+1}(k + 1 - t)\dd t =
\frac12$: hence $\frac{1}{2(k+1)^2} \leq w_k \leq
\frac{1}{2k(k+1)}$.

\textbf{4.} Upper: $\sum_{k \geq n} \frac{1}{2k(k+1)} =
\frac12\sum_{k\geq n}\bigl(\frac1k - \frac{1}{k+1}\bigr) =
\frac{1}{2n}$ (telescoping). Lower: $\frac{1}{2(k+1)^2} \geq
\frac{1}{2(k+1)(k+2)}$, whose sum telescopes to
$\frac{1}{2(n+1)}$. With question 3:
\[
\frac{1}{2(n+1)} \leq H_n - \ln n - \gamma \leq \frac{1}{2n} .
\]

\textbf{5.} Subtract $\frac{1}{2n}$: $\gamma_n - \gamma \in
\intcc{\frac{1}{2(n+1)} - \frac{1}{2n}}{0} =
\intcc{-\frac{1}{2n(n+1)}}{0}$: the corrected estimate is exact
to $O\bigl(\frac{1}{n^2}\bigr)$, and always from below.

\textbf{6.} $\gamma_{100} = 5.1873775 - \ln 100 - 0.005 =
0.5772073$, with $0 \leq \gamma - \gamma_{100} \leq
\frac{1}{20200} < 5\cdot10^{-5}$: hence $0.577207 \leq \gamma
\leq 0.577257$, i.e.\ $\gamma = 0.5772 \pm 5\cdot10^{-5}$ (true
value $0.5772156\dots$) --- four certified decimals from a
hundred terms, against ten thousand for question 2.

\textbf{7.} First dividend:
\[
H_{2m} - H_m = \Bigl(\ln 2m + \gamma + \frac{1}{4m}\Bigr)
- \Bigl(\ln m + \gamma + \frac{1}{2m}\Bigr) +
O\Bigl(\frac{1}{m^2}\Bigr)
= \ln 2 - \frac{1}{4m} + O\Bigl(\frac{1}{m^2}\Bigr).
\]
Second: the even reciprocals up to $2m$ sum to $\frac12 H_m$, so
$\sum_{j=1}^{m}\frac{1}{2j-1} = H_{2m} - \frac12 H_m = \frac12
\ln m + \ln 2 + \frac\gamma2 + o(1)$, and $\sum_{j=1}^m
\frac{1}{2j} = \frac12\ln m + \frac\gamma2 + o(1)$.

\textbf{8.} Splitting off the even terms twice:
$\sum_{k=1}^{2m}\frac{(-1)^{k-1}}{k} = H_{2m} - 2\cdot\frac12
H_m = H_{2m} - H_m$. By question 7 this equals $\ln 2 -
\frac{1}{4m} + O(m^{-2})$: the sum is $\ln 2$
(\cref{ex:b1:series:ln2} again) \emph{with} its speed.

\textbf{9.} For $N = 2m$: $S - S_N = \frac{1}{4m} + O(m^{-2}) =
\frac{1}{2N} + O(N^{-2})$. For $N = 2m + 1$: $S_{N} = S_{2m} +
\frac{1}{2m+1}$, so
\[
S - S_N = \Bigl(\frac{1}{4m} - \frac{1}{2m+1}\Bigr) +
O\Bigl(\frac{1}{m^2}\Bigr)
= -\frac{1}{4m} + O\Bigl(\frac{1}{m^2}\Bigr)
= -\frac{1}{2N} + O\Bigl(\frac{1}{N^2}\Bigr).
\]
In both cases $S - S_N \sim \frac{(-1)^N}{2N}$: half the
worst-case bound $a_{N+1}$, with a known, alternating sign.

\textbf{10.} Averaging kills the oscillating leading term:
\[
S - \tilde S_N = \frac{(S - S_N) + (S - S_{N+1})}{2}
= \frac{(-1)^N}{2}\Bigl(\frac{1}{2N} - \frac{1}{2(N+1)}\Bigr)
+ O\Bigl(\frac{1}{N^2}\Bigr) = O\Bigl(\frac{1}{N^2}\Bigr).
\]
Numerically: $S_{10} = 0.645635$, $S_{11} = 0.736544$, $\tilde
S_{10} = 0.691089$, and $\ln 2 = 0.693147$: the error drops from
$4.8\cdot10^{-2}$ to $2.1\cdot10^{-3}$ --- one addition, twenty
times better.

\textbf{11.} The alternating error changes sign at every step,
so consecutive partial sums straddle the limit and their mean
cancels the first-order term; the bracket error $H_n - \ln n -
\gamma \approx \frac{1}{2n}$ has constant sign, so no averaging
along $n$ can cancel it. Averaging the two \emph{brackets} does
help: $\frac{a_n + b_n}{2} = H_n - \ln\sqrt{n(n+1)}$, and since
$\ln\sqrt{n(n+1)} = \ln\bigl(n + \frac12\bigr) + O(n^{-2})$,
\[
H_n - \ln\Bigl(n + \frac12\Bigr)
= \Bigl(H_n - \ln n - \frac{1}{2n}\Bigr) + \frac{1}{8n^2} +
O\Bigl(\frac{1}{n^3}\Bigr) = \gamma + O\Bigl(\frac{1}{n^2}\Bigr)
\]
(question 5 and $\ln(1 + \frac{1}{2n}) = \frac{1}{2n} -
\frac{1}{8n^2} + O(n^{-3})$). Check at $n = 10$: $H_{10} -
\ln 10.5 = 0.57759$, already within $4\cdot10^{-4}$ of
$\gamma$.

\textbf{12.} $\sum_{j\leq m} \frac{1}{2j-1} \geq \sum_{j \leq m}
\frac{1}{2j} = \frac12 H_m \to +\infty$: both the positive and
the negative part of the alternating harmonic series diverge.

\textbf{13.} Write $u_n^\pm = \frac{\abs{u_n} \pm u_n}{2} \geq
0$, so $u_n = u_n^+ - u_n^-$ and $\abs{u_n} = u_n^+ + u_n^-$. If
$\sum u_n^+$ converged, then $\sum u_n^- = \sum (u_n^+ - u_n)$
would converge (difference of convergent series), hence $\sum
\abs{u_n}$ too: contradiction with conditional convergence. By
symmetry both $\sum u_n^\pm$ diverge (to $+\infty$): an infinite
reservoir of positive and of negative mass.

\textbf{14.} Let $S = \sum u_n$, $\varepsilon > 0$, and $N$ with
$\sum_{n > N} \abs{u_n} \leq \varepsilon$ (Cauchy criterion for
$\sum\abs{u_n}$). Let $M_0$ be large enough that $\sigma(\{0,
\dots, M_0\}) \supseteq \{0, \dots, N\}$. For $M \geq M_0$, the
difference $\sum_{m \leq M} u_{\sigma(m)} - \sum_{n \leq N} u_n$
is a finite sum of \emph{distinct} terms $u_n$ with $n > N$,
hence of absolute value $\leq \varepsilon$; and $\abs{S -
\sum_{n\leq N} u_n} \leq \varepsilon$ as well. So the rearranged
partial sums are within $2\varepsilon$ of $S$ eventually:
$\sum u_{\sigma(n)} = S$. Absolute convergence is
rearrangement-proof.

\textbf{15.} Each phase of the greedy procedure ends after
finitely many terms, because the remaining positive
(respectively negative) terms alone have divergent partial sums
(question 12): the running sum must eventually cross $t$. The
procedure therefore alternates infinitely many finite phases,
consuming the positive terms in order and the negative terms in
order: every term is used exactly once --- a rearrangement.
After the first crossing, between two consecutive crossings the
partial sums move monotonically toward $t$, and at a crossing
they overshoot by at most the term just added; since the terms
used at the $j$-th crossing have index at least $j$ in their
class, these overshoots tend to $0$. Hence the partial sums
converge to $t$: every real number is the sum of some
rearrangement.

\textbf{16.} The positive slots receive $\frac{1}{2j-1}$ for
$j = 1, 2, \dots$ in order, the negative slots $\frac{1}{2j}$ in
order: every term of the alternating harmonic series appears
exactly once. For $(p, q) = (1, 2)$, the blocks are
$\bigl(1, -\frac12, -\frac14\bigr)$, $\bigl(\frac13, -\frac16,
-\frac18\bigr)$, $\bigl(\frac15, -\frac1{10},
-\frac1{12}\bigr)$, \dots\ --- the displayed series.

\textbf{17.} Since $\frac{1}{4k-2} = \frac12\cdot\frac{1}{2k-1}$:
\[
\frac{1}{2k-1} - \frac{1}{4k-2} - \frac{1}{4k}
= \frac12\,\frac{1}{2k-1} - \frac12\,\frac{1}{2k}
= \frac12\Bigl(\frac{1}{2k-1} - \frac{1}{2k}\Bigr).
\]
Summing over $k = 1, \dots, K$: $T_{3K} = \frac12
\sum_{k=1}^{K}\bigl(\frac{1}{2k-1} - \frac{1}{2k}\bigr) =
\frac12 S_{2K}$: at every third partial sum, the rearranged
series is \emph{exactly} half the original one.

\textbf{18.} After $K$ complete blocks, the rearranged partial
sum is $\sum_{j=1}^{pK}\frac{1}{2j-1} -
\sum_{j=1}^{qK}\frac{1}{2j}$, and question 7 evaluates it:
\[
\Bigl(\frac{\ln(pK)}{2} + \ln 2 + \frac\gamma2\Bigr)
- \Bigl(\frac{\ln(qK)}{2} + \frac\gamma2\Bigr) + o(1)
= \ln 2 + \frac12\ln\frac pq + o(1) :
\]
the $\gamma$s cancel, the $\ln K$s cancel, the ratio
$\frac pq$ survives.

\textbf{19.} A partial sum inside block $K + 1$ differs from the
$K$-block sum by at most $p + q$ terms, each of absolute value
$\leq \frac{1}{2qK}$-ish, hence by $O\bigl(\frac1K\bigr) \to 0$:
the full sequence of partial sums has the same limit $\ln 2 +
\frac12\ln\frac pq$. For $(1, 2)$: $\ln 2 + \frac12\ln\frac12 =
\frac{\ln 2}{2} = 0.34657\dots$, and indeed $T_9 = 0.30833$
creeps toward it: by question 17, $T_{3K} = \frac12 S_{2K}$
converges with exactly half the alternating-harmonic error.
Same terms, half the sum.

\textbf{20.} $(1,1)$: $\ln 2 + \frac12\ln 1 = \ln 2$ --- the
original order, consistency. $(2,1)$: $\frac32\ln 2 \approx
1.0397$. The $(p,q)$ menu reaches exactly the countable dense
family $\ln 2 + \frac12\ln r$, $r \in \Q_{>0}$; Riemann's greedy
recipe (question 15) reaches \emph{every} real. Structure buys
formulas; greed buys totality.

\textbf{21.} The partial sums telescope: $\sum_{k=1}^{N}
\bigl(\frac1k - \ln\frac{k+1}{k}\bigr) = H_N - \ln(N+1) = b_N
\to \gamma$: the series of \cref{exo:b1:series:7} sums exactly
to Euler's constant.

\textbf{22.} Greedy for $t = 1$: the first positive term brings
the sum exactly to $1$, not beyond it, so a second positive is
taken to cross: $1, \frac13$ (sum $1.3333 > 1$), then
$-\frac12$ ($0.8333$), $\frac15$ ($1.0333$), $-\frac14$
($0.7833$), $\frac17, \frac19$ ($1.0373$), $-\frac16$
($0.8706$), $\frac1{11}, \frac1{13}$ ($1.0384$), $-\frac18$
($0.9134$), $\frac1{15}$ ($0.9801$), \dots\ --- the sums breathe
around $1$ with ever smaller amplitude, two positives now
needed per cycle since the negatives are larger.

\textbf{23.} Build blocks: at stage $j$, append enough unused
positive terms to raise the partial sum by at least $1$
(possible: the remaining positive terms have divergent sums,
question 12), then append the single negative term $-\frac{1}
{2j}$. Every positive term is eventually used (each stage uses
at least one), every negative one too (one per stage): a
rearrangement. Each stage changes the sum by $\geq 1 -
\frac{1}{2j} \geq \frac12$: the partial sums exceed
$\frac{j}{2}$ after stage $j$ and the increments within a stage
are positive except the last, bounded by $\frac{1}{2j} \to 0$:
divergence to $+\infty$.

\textbf{24.} By the law, $H_N \geq 20 \iff \ln N \geq 20 -
\gamma - \frac{1}{2N} + O(N^{-2})$: the threshold $N^*$
satisfies $\ln N^* = 20 - \gamma + o(1)$, i.e.\ $N^* =
\eu^{20 - \gamma}(1 + o(1)) \approx \eu^{19.4228} \approx
2.72\cdot10^{8}$ --- inside the crude window
$\intcc{1.8\cdot10^8}{4.9\cdot10^8}$ of
\cref{ex:b1:series:harmonicstack}, and pinned by $\gamma$.

\textbf{25.} (i) In $H_n = \ln n + \gamma + \frac{1}{2n} +
O(n^{-2})$: the $\ln n$ is the integral, $\gamma$ the price of
replacing a sum by an integral (a genuinely new constant of
analysis), and $\frac{1}{2n}$ the first correction --- the
trapezoid's shadow. (ii) Conditional convergence leans on
cancellation between two infinite reservoirs (question 13), so
reordering re-weights the reservoirs; absolute convergence has
finite total mass, and question 14's tail estimate is
order-blind. (iii) The $(p,q)$ sums were \emph{computed}: the
$\gamma$-law turned each rearranged partial sum into $\ln 2 +
\frac12\ln\frac pq + o(1)$, with $\gamma$ itself canceling ---
an asymptotic bookkeeping exercise, not an abstract argument.
(iv) Next: products and unconditional summability for power
series in the Year 2 volume; and the Year 3 volume's weekend
problem on Stirling's formula, where sum-versus-integral
bookkeeping, pushed one order further, produces $\sqrt{2\pi}$
itself.
\end{solution}
